
\section{Middle primes III: the reference law $\mu$, first stage $\lambda$, and its dominator}
\label{sec:middle-frame}

\paragraph{What this section proves and how to read it.}
This section builds $\mu$ (five-type completion of the
twelve-label principal law), $\lambda$ (polarization and one-stage
replacements), and the law $\Nbr$ dominating $\lambda$ at $\iota$,
and proves (CP1)--(CP3) for middle atoms with prices $L_0,B$.
Inputs: caps of \cref{sec:middle-reference,sec:positive-roots}.
Outputs: $\mu,\lambda,\Nbr$ for \cref{sec:middle-records}. A first
reading should follow the three constructions and the (CP1)--(CP3)
statements; weight verification is reference material.
``Frame'' in labels below means the package of these three laws together with the records and prices defined with them, as
defined in \cref{def:record:frame}. On a middle block the
reference (\cref{subsec:middle-reference-law}) is the five-type
completion of the twelve-label principal law, which is uniform on
the intervals below $2^{30}$ and tilted
(\cref{sec:positive-roots}) on the seven intervals from $2^{30}$;
its completion weights use the caps of the principal law
(\cref{sec:middle-reference,sec:positive-roots}) and of the
reference fold record (\cref{lem:old-f-law}). The first stage
(\cref{subsec:middle-first-stage}) is obtained by polarization at
one chamber cut, which removes the two nonaccepting one-layer
fold copies $\operatorname{Diag}_\iota(\Pi^{(1)}_{F_1})$,
$\operatorname{Diag}_\iota(\Pi^{(1)}_{F_2})$, and by replacing the
one-layer matching copies by accepting laws with the same actual
conditional marginals; the two resulting record differences are
the tags of (CP2)(b). The law $\Nbr$
(\cref{subsec:middle-gram}) then dominates $\lambda$ at $\iota$ in
a $2\times2$ block, giving (CP3), and
\cref{subsec:middle-frame-prices} computes the resulting prices
$L_0$ and $B$. Later sections build the comparison records on
this reference (\cref{sec:middle-records}) and compare them with
$\mu$, $\lambda$ and $\Nbr$ through the variance savings
(\cref{sec:variance}), the dominating matrices
(\cref{sec:envelopes}) and the matching envelopes
(\cref{sec:matching-mean}); the closing
(\cref{sec:closing}) uses the prices through $B$, $R_{\rm err}$
and $L_0$.

The middle reference uses twelve labels, two layers over one chamber
coordinate, and is pulled back to the slot array $V$ by ignoring the second
chamber coordinate. It is the five-type completion of \cref{lem:five-type}
applied to the twelve-label root: folds, matchings, merges and constants are
added with weights in powers of $r=b^{-1/2}$ ($b$ the least prime of the block) until every chamber cut is
positive, and one-layer copies placed on the reserved diagonal restore
positivity at the layer flip $\iota$. On the seven tilted pieces the positive
root changes two inputs of the completion, its caps at the chamber cut L20 and
at the layer cut L21, both proved in fixed half charts
(\cref{rem:root-charts}). (\Cref{sec:history} explains how this
completion developed from an earlier version on $216$ slots;
nothing here depends on that history.) The comparison law built
on this reference need not be reflection positive; it is compared with
the first stage $\lambda$ only at $\iota$, through the
$2\times2$ block of \cref{subsec:middle-gram}, with symmetric and
skew errors priced separately.

\emph{Conventions.} In this section $\vartheta_a$ is the chamber cut
$\vartheta_{0,a}$ of the first coordinate: the reflection
$u_0\mapsto2a-1-u_0$ with plus half $H_a=\{a,a+1,a+2\}$ (indices mod $6$),
$a\in\Z/6\Z$. The reserved cut $\iota$ flips the layer $\ell$ and has plus
half $V_+=\{\ell=0\}$. For a prime $p$ we write $r=p^{-1/2}$ and $q=1/p$; on a
block $\Blk$ with least prime $b=b_{\Blk}$, $r=b^{-1/2}$ and $q=1/b$, while
$q_{\Blk}=\prod_{p\in\Blk}p$. The order $\succeq_c$ of cut matrices is that of
\cref{def:cc:order}. The pieces, the blocks, the twelve-label graphs and the
caps of their cuts are those of \cref{sec:middle-reference,sec:positive-roots}.
Every middle block has $\log q_{\Blk}\le L=570$; the local data of a piece are
evaluated at its start $X$ and at the local cap $L_{\rm loc}=585$
(\cref{lem:record:local-domain}).

\subsection{The reference law of a middle block}
\label{subsec:middle-reference-law}

Fix a middle block $\Blk$ of a piece starting at $X$, with least prime $b$,
$r=b^{-1/2}$ and $q=1/b$. An \emph{object} carries one label of $X_{\Blk}$
per layer. Write $\lambda_{\rm F}=\lambda_{\Blk}(G_{\rm F,2})$ and
$\lambda_{\rm E}=\lambda_{\Blk}(G_{\rm E,2})$ for the two-layer $F$ and $E$
laws of three objects $(y_0,y_1,y_2)$, $y_o=(y_{o,0},y_{o,1})$, and
$\lambda^{(1)}_{\rm F}$, $\lambda^{(1)}_{\rm E}$ for their one-layer versions.
The components of the reference are images of these laws under the patterns
of \cref{def:patterns}: the folds $F_0=(0,1,2,2,1,0)$, $F_1=(1,0,0,1,2,2)$ and
$F_2=(2,2,1,0,0,1)$ applied to the $F$ law, and the matchings
$E_0=(0,1,1,2,2,0)$ and $E_1=(0,0,1,1,2,2)$, the six merges $G_{sk}$ (which use
the objects $0,1$) and the two constants $C_0,C_1$ (which use the object $0$)
applied to the $E$ law. Under a pattern $\phi$, a configuration
$(y_0,y_1,y_2)$ of three objects gives the slot $(u_0,u_1,\ell)$ of $V$ the
label $y_{\phi(u_0),\ell}$, whatever $u_1$ is. We write $\Pi_\phi$
($\Pi_{\phi,\Blk}$ when the block matters) for the resulting image on
$X_{\Blk}^V$ of the two-layer law, $\Pi^{(1)}_\phi$ for that of the one-layer
law, and $\Pi_{\rm L}$ for the two-layer root of the block pulled back in the
same way: the uniform law $\lambda_{\Blk}(G_{\rm L,2})$ on the pieces below $2^{30}$, and on the pieces from $2^{30}$ the positive root
$\Gamma_{\Blk,\ulaw}$ of \cref{def:positive-root} with the parameter of the
piece; $\Pi^{(1)}_{\rm L}$ is the pullback of the one-layer root
$\lambda_{\Blk}(G_{\rm L,1})$, which is not tilted. The five types
$\rm L,D,S,R,C$, with multiplicities $1,3,2,6,2$ and degrees $0,1,1,2,3$, are
those of \cref{def:patterns}; $\pi_{\mathfrak t,j}$ is the sum of the $j$-layer
components of type $\mathfrak t$, of mass equal to the multiplicity. For a law
$\kappa$ of one-layer configurations, $\operatorname{Diag}_\iota(\kappa)$
places the same configuration on both layers (the diagonal copy of
\cref{lem:commoncut}(i)); it is fixed samplewise by $\iota$. The halves of the
components at $\vartheta_0$ are listed in \cref{tab:corner}.

\paragraph{Rows.}
Let $j\in\{1,2\}$. Under the $j$-layer $E$ law let $\nu$ be the law of one object (in this paragraph $\nu$ is this one-object law, not the comparison law)
and $\nu^\xi$ the conditional law of an outer object given that the
middle object $y_1$ equals $\xi$: $\nu$ is uniform on the $j$-tuples, indexed by the layers, of block
labels whose labels at each prime are pairwise distinct, and $\nu^\xi$ is uniform on those whose labels
avoid the labels of $\xi$ at every prime. Let $T^\xi_{\rm F}$ and
$T^\xi_{\rm E}$ be the densities, relative to $\nu^\xi\otimes\nu^\xi$, of the
conditional law of $(y_0,y_2)$ given $y_1=\xi$ under the $j$-layer $F$ and
$E$ laws, regarded as operators on $L^2(\nu^\xi)$. For a prime $p$ let
$\rho_{{\rm F},j}(p)$ and $\rho_{{\rm E},j}(p)$ be the Schur-test bounds of
\cref{cor:influence-schur} for the conditional maximal correlation of the
labels of objects $0$ and $2$, with the labels of object $1$ pinned, under
$\lambda_p(G_{{\rm F},j})$ and $\lambda_p(G_{{\rm E},j})$, in the form of
\cref{rem:influence-uses}(1) (graded count errors; the two vertex sets of the
corollary are the labels of the two outer objects); they are explicit
functions of $r=p^{-1/2}$, and
$\rho_{{\rm F},j}(X)$, $R_j(X)$ below denote their values at $r=X^{-1/2}$.
With $n=L_{\rm loc}/\log X$,
$R_1(p)=(p-1)/(p-3)$, $R_2(p)=(p-3)(p-7)/((p-11)(p-15))$ and
$\upsilon_j=(e^{n(R_j(X)-1)^2}-1)^{1/2}$, the \emph{row constants} of the
piece are
\[
 f_j=\frac{R_j(X)^n\rho_{{\rm F},j}(X)+2\upsilon_j+\upsilon_j^2}{X^{-1/2}},
 \qquad h_j=\frac{\rho_{{\rm E},j}(X)}{X^{-1/2}} .
\]
These are the row constants that the certificate evaluates
(\cref{rem:row-comparison-uses}(1)).

\begin{lemma}[The reference]\label{lem:front12:reference}
Let $\Blk$ be a middle block of a piece starting at $X$, with least prime $b$
and $r=b^{-1/2}$.
\begin{enumerate}[label=(\roman*)]
\item \emph{Rows.} For $j\in\{1,2\}$ the middle object has the law $\nu$
 under the $j$-layer $F$ law as well as under the $j$-layer $E$ law, and
 for every value $\xi$,
 \[
  \|T^\xi_{\rm F}-\Jproj\|\le f_jr,\qquad \|T^\xi_{\rm E}-\Jproj\|\le h_jr
 \]
 on $L^2(\nu^\xi)$, $\Jproj$ being the projection onto constants.
\item \emph{Negative envelopes.} Let $d$, $d_1$, $k_{\rm L}$ and $k_{\rm E}$
 be the block caps, at the piece start $X$ and the local cap $L_{\rm loc}$, of
 the root cuts L20 and L10 and of the layer cuts L21 and E21: for L10 and
 E21, and for L21 on the pieces below $2^{30}$, those of
 \cref{lem:layer12-caps}(b); for L20 on the pieces below $2^{30}$ that of
 \cref{imp:local:collision-gap}; for L20 and L21 on the pieces from $2^{30}$ those
 of the positive root (\cref{cert:positive-root-caps}). Let $k_{\rm D}=k_{\rm D}(X,L_{\rm loc})$ be
 the cap of the reference fold record of \cref{lem:old-f-law}(e), and put
 $k_{\rm S}=k_{\rm R}=k_{\rm C}=k_{\rm E}$. Then
 \[
  \Pi_{\rm L}\succeq_{\vartheta_0}-dr\,\Pi_{F_0},\qquad
  \Pi^{(1)}_{\rm L}\succeq_{\vartheta_0}-d_1r\,\Pi^{(1)}_{F_0},
 \]
 and $\Pi_\phi\succeq_\iota-k_{\mathfrak t}r\operatorname{Diag}_\iota(\Pi^{(1)}_\phi)$
 for every two-layer component $\Pi_\phi$ of type $\mathfrak t$, the root
 included.
\item \emph{Completion.} Let $w=(w_{\rm L},w_{\rm D},w_{\rm S},w_{\rm R},w_{\rm C})$
 with $w_{\rm L}=1$ and $v=(v_{\rm L},v_{\rm D},v_{\rm S},v_{\rm R},v_{\rm C})$ be
 nonnegative weights satisfying the thirteen inequalities
 \begin{align*}
  &w_{\rm D}\ge d,&&w_{\rm S}\ge w_{\rm D},&&w_{\rm R}\ge f_2w_{\rm D}+h_2w_{\rm S},
   &&w_{\rm C}w_{\rm S}\ge w_{\rm R}^2,\\
  &v_{\rm D}\ge d_1v_{\rm L},&&v_{\rm S}\ge v_{\rm D},&&v_{\rm R}\ge f_1v_{\rm D}+h_1v_{\rm S},
   &&v_{\rm C}v_{\rm S}\ge v_{\rm R}^2,\\
  &v_{\mathfrak t}\ge k_{\mathfrak t}w_{\mathfrak t}\quad(\mathfrak t\in\{\rm L,D,S,R,C\}).
 \end{align*}
 Then
 \[
  \widetilde\mu_{\Blk}=\sum_{\mathfrak t}w_{\mathfrak t}r^{\deg\mathfrak t}\pi_{\mathfrak t,2}
   +r\sum_{\mathfrak t}v_{\mathfrak t}r^{\deg\mathfrak t}\operatorname{Diag}_\iota(\pi_{\mathfrak t,1})
 \]
 is a nonnegative measure on $X_{\Blk}^V$ of mass at least one, reflection
 positive at every cut of the family (the chamber cuts of both
 coordinates and $\iota$), and invariant under the layer swap and under the
 square-affine maps of (R3) of \cref{hyp:record:reference}; in particular
 its one-slot marginals are its mass times the uniform law.
\item \emph{Protection.} If the weights satisfy moreover the fourteenth
 inequality
 \begin{equation}\label{eq:protection}
  v_{\rm D}\ge2k_{\rm D}w_{\rm D},
 \end{equation}
 which implies $v_{\rm D}\ge k_{\rm D}w_{\rm D}$ of the third line, then for
 $s\in\{1,2\}$
 \[
  2w_{\rm D}r\,\Pi_{F_s}+v_{\rm D}r^2\operatorname{Diag}_\iota(\Pi^{(1)}_{F_s})
  \succeq_\iota0.
 \]
 Parts (i)--(iii) do not use \eqref{eq:protection}; the bracket of
 \cref{subsec:middle-gram} does.
\end{enumerate}
The \emph{reference law} of $\Blk$ is
$\mu_{\Blk}=\widetilde\mu_{\Blk}/\widetilde\mu_{\Blk}(1)$. It satisfies (R1)--(R3)
of \cref{hyp:record:reference}, with (R2) at every cut of the family, and it
is admissible (\cref{def:admissible}).
\end{lemma}

\begin{proof}
(i) \emph{One prime.} For one layer the middle law is uniform for both laws,
by translation invariance; for two layers it is uniform on the distinct
pairs for the $E$ law by symmetry and for the $F$ law by
\cref{lem:record:class-three}(ii), since $p\ge Q_3>11$. Under the $E$ law the
outer objects given $y_1=\xi$ have exactly the marginal $\nu^\xi$, so
$T^\xi_{\rm E}$ is the density of a coupling with exact marginals and
$\|T^\xi_{\rm E}-\Jproj\|$ is its conditional maximal correlation, at most
$\rho_{{\rm E},j}(p)$ by \cref{cor:influence-schur}. Under the $F$ law the
actual outer marginals given $y_1=\xi$ have densities in
$[R_j(p)^{-1},R_j(p)]$ relative to $\nu^\xi$ (\cref{lem:record:class-three}(iii),
$p>15$), and the coupling relative to these actual marginals has maximal
correlation at most $\rho_{{\rm F},j}(p)$; in the language of the four-count
comparison (\cref{lem:four-count}), these densities are ratios of completion
counts.

\emph{The block.} Fix $j\in\{1,2\}$. Given $\xi=(\xi_p)_p$ the conditional laws are products over
the primes of $\Blk$. By \cref{lem:row-comparison}(c), the densities of the
actual outer marginals multiply, so they are at most
$R_{\Blk}=\prod_pR_j(p)\le R_j(b)^n$, because $R_j$ decreases and
$|\Blk|\le n_X\le n$ (\cref{lem:record:local-domain}(ii)); their chi-square
defects multiply, and a density with values in $[R^{-1},R]$ has defect at
most $R-1$; writing $g$ for an outer marginal density and $g_p$ for its factor at $p$,
\[
 1+\|g-1\|^2=\prod_p(1+\|g_p-1\|^2)\le\exp\Bigl(\sum_p(R_j(p)-1)^2\Bigr)
 \le e^{n(R_j(b)-1)^2}
\]
for each such $g$; and the maximal correlation of the
product coupling is the largest of the factors, at most
$\max_p\rho_{{\rm F},j}(p)$, respectively $\max_p\rho_{{\rm E},j}(p)$. Hence
\cref{lem:row-comparison}(a), with $H_{\rm L},H_{\rm R}\le R_{\Blk}$ and
$v_{\rm L},v_{\rm R}\le\upsilon$, $\upsilon^2=e^{n(R_j(b)-1)^2}-1$, gives
$\|T^\xi_{\rm F}-\Jproj\|\le R_{\Blk}\max_p\rho_{{\rm F},j}(p)+2\upsilon+\upsilon^2$;
the $E$ kernel has exact marginals, so
$\|T^\xi_{\rm E}-\Jproj\|\le\max_p\rho_{{\rm E},j}(p)$.

\emph{Evaluation at the piece start.} As functions of $r=p^{-1/2}$ (with $n$
fixed), $(R_1-1)/r^2=2/(1-3r^2)$ and
$(R_2-1)/r^2=(16-144r^2)/((1-11r^2)(1-15r^2))
=\sum_{m\ge0}(24\cdot15^m-8\cdot11^m)r^{2m}$ are power series with
nonnegative coefficients, so $\upsilon/r$ and $\upsilon^2/r=r(\upsilon/r)^2$ are
admissible majorant functions in the sense of \cref{lem:majorant} (a different use of \emph{admissible} from admissible laws), since $(e^x-1)/x$ increases; $R_j^n$ is
nondecreasing in $r$; and $\rho_{{\rm F},j}/r$, $\rho_{{\rm E},j}/r$ are admissible,
the bound of \cref{cor:influence-schur} being $r$ times an admissible function
of $r$ (\cref{rem:influence-uses}(3)). Hence $\max_p\rho(p)\le\rho(b)$, and the bound
of the previous paragraph, divided by $r$, is a majorant function of $r$,
so the bound is at most $r$ times that quotient at $r=b^{-1/2}$, which is at most its value at $X$ by \cref{lem:majorant}(b):
$\|T^\xi_{\rm F}-\Jproj\|\le f_jr$ and $\|T^\xi_{\rm E}-\Jproj\|\le h_jr$.

(ii) The plus half of $\vartheta_0$ carries the objects $0,1,2$ of $H_0$, and
by \cref{lem:layer12-geometry}(b) the half graph of the root at $\vartheta_0$
is $G_{\rm F,2}$, whose labelling law $\lambda_{\rm F}$ is the law of the plus
half of $\Pi_{F_0}$. The configurations of $\Pi_{F_0}$ are fixed by
$\vartheta_0$ (as $F_0(5-u)=F_0(u)$), so the cut form of $\Pi_{F_0}$ at a real
half-function $\Phi$ is $\E_{\lambda_{\rm F}}\Phi^2$. For the uniform root the
cut form of $\Pi_{\rm L}$ is $\ip\Phi{\mathsf A_{\Blk}\Phi}$, where $\mathsf A_{\Blk}$ is its cut operator on $L^2(\lambda_{\rm F})$, and
\cref{imp:local:collision-gap} for the cut L20 gives the first inequality. For the
positive root of a piece from $2^{30}$ the cut form of $\Pi_{\rm L}$ at $\Phi$
is $\ip\Phi{T_{\Blk}\Phi}$ on $L^2(\lambda_{\rm F})$, where $T_{\Blk}$ is the
cut operator of the reweighted graph law relative to the fixed half law
$\lambda_{\rm F}$ (\cref{lem:spec:cutkernel}(f)); \cref{cert:positive-root-caps}
with \cref{lem:spec:tensorcomp}(b) gives $T_{\Blk}\succeq-dr$, which is the
first inequality with the same $\Pi_{F_0}$ (\cref{rem:root-charts}). The
one-layer root is the uniform law on every piece, and
\cref{lem:layer12-caps}(b) for the cut L10 gives the second inequality. At
$\iota$, by \cref{lem:layer12-geometry}(c) the half graphs of $G_{\rm L,2}$ and
$G_{\rm E,2}$ are their one-layer graphs, and the cut form of
$\operatorname{Diag}_\iota(\Pi^{(1)}_\phi)$ at $\Phi$ is the mean of $\Phi^2$
under $\Pi^{(1)}_\phi$. So the caps L21 and E21 give
$\lambda_{\Blk}(G_{\rm L,2})\succeq_\iota-k_{\rm L}r\operatorname{Diag}_\iota(\lambda_{\Blk}(G_{\rm L,1}))$
for the uniform root and
$\lambda_{\rm E}\succeq_\iota-k_{\rm E}r\operatorname{Diag}_\iota(\lambda^{(1)}_{\rm E})$.
For the positive root, the layer bounds of \cref{cert:positive-root-caps} with
\cref{lem:spec:cutkernel}(f) and \cref{lem:spec:tensorcomp}(b) give the
inequality relative to the half law of
$\operatorname{Diag}_\iota(\Pi^{(1)}_{\rm L})$, the untilted one-layer root:
$\Gamma_{\Blk,\ulaw}\succeq_\iota
-k_{\rm L}r\operatorname{Diag}_\iota(\lambda_{\Blk}(G_{\rm L,1}))$.
A component is the image of its field law under an emission, and it inherits
the bound of the field law (\cref{lem:five-type}(e)); this gives the
inequality for the root and, with $k_{\rm E}$, for the matchings, the merges
and the constants. For the folds it is \cref{lem:old-f-law}(e), with the cap
$k_{\rm D}=k_{\rm D}(X,L_{\rm loc})$: its values at the $33$ piece starts,
with the gates of that lemma, are in \cref{cert:record:old-f-cap}, and its
conditions on the block ($p\ge X\ge2^{24}$,
$\sum_{p\in\Blk}\log p\le L_{\rm loc}$) hold by
\cref{lem:record:local-domain}(i),(ii).

(iii) \emph{Positivity, mass and invariance.} All components are probability
laws and all weights are nonnegative; the root has weight $w_{\rm L}=1$, so
$\widetilde\mu_{\Blk}(1)\ge1$. Every graph law $\lambda_p(G)$ is invariant under
the diagonal maps $z\mapsto\alpha z+\beta$ with $\alpha$ a nonzero square,
which preserve injectivity and multiply differences by squares; the
component maps act on the vertices only and commute with them, and on a
block the maps act independently at each prime. The positive root is
invariant under the same maps and under the layer swap
(\cref{lem:positive-root-properties,lem:imp:coherent39-law}). Translations act transitively, so every
one-slot marginal is its mass times the uniform law (\cref{lem:marginals}). The layer swap is an
automorphism of the three two-layer graphs (\cref{lem:layer12-geometry}(a)),
and a diagonal copy is fixed by it.

\emph{First-coordinate cuts.} For $j\in\{1,2\}$, the middle law $\nu$ and the
conditional references $\nu^\xi$ of the $j$-layer $E$ law, the two-object marginal of that law, the $j$-layer $E$ and $F$ laws and the $j$-layer root form a
five-type datum on the label set of one object (\cref{def:five-type-datum}):
the $E$ law is exchangeable; by (i) the middle law of the $F$ law is $\nu$, so
that the $F$ law has the form required there with the kernels $T^\xi_{\rm F}$;
and the root is invariant under $\mathcal G_6$ (\cref{lem:layer12-geometry}(a)
for the uniform root, \cref{lem:positive-root-properties,lem:imp:coherent39-law}
for the positive root). For $j=2$, (i) gives the full-norm hypotheses
of \cref{lem:five-type}(b) with $e_{\rm F}=f_2r$ and $e_{\rm E}=h_2r$, and (ii)
the root envelope with $\delta=dr$; for the graded weights
$(1,w_{\rm D}r,w_{\rm S}r,w_{\rm R}r^2,w_{\rm C}r^3)$ the conditions of
\cref{lem:five-type}(b) are the four inequalities of the first line. So
$\sum_{\mathfrak t}w_{\mathfrak t}r^{\deg\mathfrak t}\pi_{\mathfrak t,2}$ is
reflection positive at every chamber cut of the first coordinate. The same
holds for $j=1$ with the inputs $d_1,f_1,h_1$ and the weights $v$ (the second
line, homogeneous in $v_{\rm L}$), and the diagonal copy preserves this
positivity (\cref{lem:five-type}(e)).

\emph{The reserved cut.} By (ii) every two-layer component of type $\mathfrak t$
is bounded below at $\iota$ by $-k_{\mathfrak t}r$ times the diagonal copy of its
one-layer component, and the third line is
$v_{\mathfrak t}\ge k_{\mathfrak t}w_{\mathfrak t}$; so
$\widetilde\mu_{\Blk}\succeq_\iota0$ by the fold pairing of
\cref{lem:five-type}(e). The one-layer component of a type need not be the
one-layer marginal of its two-layer component, only the order of (ii) being
used (\cref{lem:five-type}(e)); for the positive root it is the untilted
$\Pi^{(1)}_{\rm L}$. The first coordinate and
the layer are the active coordinates of the pullback to $V$, so both
conclusions hold on $V$ (\cref{lem:slot-array}(v)).

\emph{Second-coordinate cuts.} $\widetilde\mu_{\Blk}$ is the pullback of a law on
the labels indexed by $(u_0,\ell)$, so every configuration is fixed by every
chamber cut of the second coordinate, and the cut matrix there is
nonnegative diagonal (\cref{lem:slot-array}(v), \cref{lem:commoncut}).

\emph{Normalization} preserves positivity and the marginal assertions, and
(R1) holds because the block and its law do not depend on the scale. The
reference is a nonnegative law depending only on the labels in the finite fields at the primes of the block, invariant under
the maps of (R3); so it has uniform whole-slot marginals and does not depend
on the interval, residue or inclusion indices, which is (O1) of
\cref{def:admissible}; no sharp bound is claimed for it, so clause (O2) is vacuous for it.

(iv) By (ii) for the type D,
$2w_{\rm D}r\,\Pi_{F_s}\succeq_\iota-2k_{\rm D}w_{\rm D}r^2
\operatorname{Diag}_\iota(\Pi^{(1)}_{F_s})$, and
$\operatorname{Diag}_\iota(\Pi^{(1)}_{F_s})\succeq_\iota0$, being fixed
samplewise by $\iota$ (\cref{lem:commoncut}); so \eqref{eq:protection} gives
the claim.
\end{proof}

\begin{remark}[The certified weights]\label{rem:middle-weights}
In each of the $33$ pieces the weights $w,v$ are fixed rationals
(\cref{tab:data:middle-weights}), chosen at $(X,L_{\rm loc})$ with the caps of
the piece, and the thirteen inequalities of \cref{lem:front12:reference}(iii)
and the protection \eqref{eq:protection} are checked exactly with the inputs
of the piece, $462$ inequalities in all (\cref{cert:record:frame-data}). In
every piece three of them fix the leading weights: $w_{\rm D}=w_{\rm S}$ is the
block cap $d$ of L20, rounded up; $v_{\rm L}$ is the block cap $k_{\rm L}$ of
L21, rounded up; and $v_{\rm D}=v_{\rm S}$ is $2k_{\rm D}w_{\rm D}$,
rounded up to the tabulated rational, so that the protection inequality holds, with equality before rounding. The inequality
$v_{\rm D}\ge d_1v_{\rm L}$ is slack, also on the seven pieces from $2^{30}$,
where $v_{\rm L}$ is the layer cap of the same positive root. Its
smallest slack on these seven pieces is greater than $5.7\cdot10^{-5}$;
the exact rational check retains this gate along with protection. The certificate checks $v_{\rm S}\ge2k_{\rm E}w_{\rm S}$, which is
stronger than required. Since every input of \cref{lem:front12:reference} is
evaluated at the piece start and is valid for every block of the piece
(\cref{lem:record:local-domain}), one set of weights serves all blocks of a
piece.
\end{remark}

\begin{remark}[The fixed half charts]\label{rem:root-charts}
The caps of the positive root are proved in two fixed charts: at the chamber
cut on $L^2(\lambda_{\rm F})$, the half law of $\Pi_{F_0}$, and at the layer
cut on $L^2(\lambda_{\Blk}(G_{\rm L,1}))$, the half law of
$\operatorname{Diag}_\iota(\Pi^{(1)}_{\rm L})$
(\cref{cert:positive-root-caps}). These are the diagonal measures
against which \cref{lem:front12:reference}(ii) is stated and which the
completion of (iii) uses, so the actual half marginal of the positive root,
whose density relative to the fixed half law differs from $1$ by $O(r)$ in the norms used below, is not needed there. Its change is paid
where the half marginal enters: in the restorations of the variance credits
and of the mean savings (\cref{lem:record:root-marginal}(iv)).
\end{remark}

Together with the laws of the other ranges (\cref{lem:frontier:iidlocal},
\cref{lem:upper-reference-rp}, \cref{lem:tail-law}), the reference laws of
the middle blocks form the reference $\mu=\bigotimes_{\Blk}\mu_{\Blk}$ of
(CP1). Of its components, the root and the aligned fold $\Pi_{F_0}$ accept
the six queries, both as two-layer components and as one-layer diagonals
(the positive root has the support of the uniform root); every other record
is removed or replaced, by the first stage below or by the comparison law
(\cref{cert:record:middle-geometry}).

\subsection{The first stage $\lambda$ (polarization and one-stage replacements)}\label{subsec:middle-first-stage}

The comparison law will not contain the two one-layer fold records
$\operatorname{Diag}_\iota(\Pi^{(1)}_{F_1})$ and
$\operatorname{Diag}_\iota(\Pi^{(1)}_{F_2})$, which do not accept the queries,
nor the folded matching records in their reference form. The first stage
removes them at the price of a second-order correction at one chamber cut,
and replaces the folded matching records by laws that accept the queries
and keep their actual conditional marginals.

\begin{definition}[The middle frame]\label{def:record:frame}
The frame consists of the following objects.
\begin{enumerate}[label=(\alph*)]
\item \emph{Reference.} $\mu=\bigotimes_{\Blk}\mu_{\Blk}$ is the reference of
 (CP1); on a middle block, $\mu_{\Blk}$ is the law of
 \cref{lem:front12:reference}.
\item \emph{Records.} Let $\Blk$ be a middle block with least prime $b$, and
 put $r=b^{-1/2}$ and $q=r^2$. Before the common normalization, define the following measures (the first two are summands of $\widetilde\mu_{\Blk}$; the next three are added in (d) below):
 \begin{itemize}
 \item $F_{\rm fold}=v_{\rm D}r^2\bigl(\operatorname{Diag}_\iota(\Pi^{(1)}_{F_1})
  +\operatorname{Diag}_\iota(\Pi^{(1)}_{F_2})\bigr)$ is the sum of the two
  nonaccepting one-layer fold components on the reserved diagonal, of mass $2v_{\rm D}r^2$;
 \item $E_{\rm fold}=v_{\rm S}r^2\bigl(\operatorname{Diag}_\iota(\Pi^{(1)}_{E_0})
  +\operatorname{Diag}_\iota(\Pi^{(1)}_{E_1})\bigr)$ is the sum of the two
  reserved-folded matching components, of mass $2v_{\rm S}r^2$;
 \item $E_{\rm add}=v_{\rm D}r^2\bigl(\operatorname{Diag}_\iota(\Pi^{(1)}_{E_0})
  +\operatorname{Diag}_\iota(\Pi^{(1)}_{E_1})\bigr)$ is a pair of reserved-folded
  matching records of total mass $2v_{\rm D}r^2$;
 \item $G=(f_1+h_1)r\,(G_e+G_o)$, with
  $G_e=v_{\rm D}r^2\operatorname{Diag}_\iota(\Pi^{(1)}_{G_{00}})$ and
  $G_o=v_{\rm D}r^2\operatorname{Diag}_\iota(\Pi^{(1)}_{G_{11}})$, is the
  \emph{polarization correction}: two merge records of common record weight
  $v_{\rm D}r^2$, of total mass $2v_{\rm D}(f_1+h_1)r^3$;
 \item $H$ adds, for each of the two two-layer matching components
  $\Pi_{E_0},\Pi_{E_1}$ of weight $w_{\rm S}r$, its actual half diagonal
  $\operatorname{Diag}_\iota(\Pi^{(1)}_{E_s})$ with coefficient
  $2w_{\rm S}r\cdot3/(b-3)$; thus $H$ has mass $12w_{\rm S}r^3/(1-3q)$.
 \end{itemize}
 All five measures are nonnegative and reserved diagonal, that is, fixed
 samplewise by $\iota$. In the decompositions below all five measures are
 divided by $\widetilde\mu_{\Blk}(1)$; only the first two are records of $\mu_{\Blk}$ itself.
\item \emph{First stage.} $E_{\rm new}$ is the nonnegative reserved-diagonal
 replacement of $E_{\rm fold}+E_{\rm add}$ given by
 \cref{lem:record:one-stage-folded-e}. It has the same mass, it accepts
 the six queries, and at its own first-coordinate cut ($\vartheta_1$ for $E_0$, $\vartheta_2$ for $E_1$) each record of
 $E_{\rm new}$ has the same actual conditional marginals as the record it
 replaces.
\item \emph{Decomposition.} Blockwise, and with all records zero outside the
 middle blocks,
 \begin{align*}
  \rho&=\mu-F_{\rm fold}+E_{\rm add}+G,&
  \lambda&=\mu-F_{\rm fold}-E_{\rm fold}+E_{\rm new}+G,\\
  \delta_E&=E_{\rm new}-E_{\rm fold}-E_{\rm add},&
  \Nbr&=\lambda+F_{\rm fold}+H,\\
  \Nbrz&=\mu+E_{\rm add}+G+H.
 \end{align*}
 Hence $\lambda=\rho+\delta_E$, $\Nbr=\Nbrz+\delta_E$,
 $\Nbr-\lambda=F_{\rm fold}+H$ and
 $\Nbr+\lambda=2\mu-F_{\rm fold}-2E_{\rm fold}+H+2E_{\rm new}+2G$; in
 particular $F_{\rm fold}$ occurs with coefficient one in $\Nbr-\lambda$. (The bracket measures $\Nbr$
 and $\Nbrz$ are unrelated to the scale $N=\Kint\Lsc$.)
\item \emph{Prices and tags.} $U_0$ is the summed price of
 $\Nbrz-\mu=E_{\rm add}+G+H$ and $P$ that of $F_{\rm fold}+E_{\rm add}+G$. The
 record differences of $\delta_E$ form two tags (\cref{def:cp2}(b)): the
 differences of type $E_0$, sharp at their first-coordinate cut
 $\vartheta_1$, form the tag $0$, and those of type $E_1$, sharp at
 $\vartheta_2$, the tag $1$. $V$ and $D$ are the summed variation and sharp
 prices of $\delta_E$, the sharp price of a difference being its mass times
 its sharp bound at its own cut, relative to its actual half marginals, times
 $q_{\Blk}^{2a}$ (\cref{def:price}(b)). The two tags have equal totals, the
 sharp total $D/2$ and the variation total $V/2$ each, because the two
 records of a block have the same mass and the same sharp bound
 (\cref{lem:record:middle-prices}). With the two-tag series $\tagser$ of
 \cref{cor:two-cut-series},
 \[
  R_{\rm err}=e^{U_0}\tagser(D/2,V/2),\qquad B=e^{U_0}+R_{\rm err},\qquad
  L_0=e^{P}\tagser(D/2,V/2).
 \]
\end{enumerate}
\end{definition}

The actual half marginal of $\Pi_{E_s}$ at $\iota$ is the uniform law on
ordered triples of distinct labels placed by $E_s$, which is the half law
of $\Pi^{(1)}_{E_s}$ (\cref{lem:record:fold-identification}(d)); this is the
diagonal that $H$ adds. All records in \cref{def:record:frame}(b) are images of invariant laws under
vertex maps, followed by a diagonal copy, so they are positive records
(\cref{def:cp:records}): their normalizations are admissible. Admissibility of the replacements in \cref{def:record:frame}(c) is verified after \cref{lem:record:one-stage-folded-e} below.

\begin{lemma}[The polarized first stage]\label{lem:middle-polarization}
For every middle block, $E_{\rm add}+G-F_{\rm fold}\succeq_{\vartheta_0}0$. Hence
$\rho_{\Blk}\succeq_{\vartheta_0}\mu_{\Blk}\succeq_{\vartheta_0}0$, and (CP2)(a)
holds at the cut $\vartheta_0$, with the variation price $p_{\Blk}$ of the
positive records $F_{\rm fold}$, $E_{\rm add}$ and $G$.
\end{lemma}

\begin{proof}
The one-layer data of the proof of \cref{lem:front12:reference}(iii) form a
five-type datum whose rows satisfy $\|T^\xi_{\rm F}-\Jproj\|\le f_1r$ and
$\|T^\xi_{\rm E}-\Jproj\|\le h_1r$ (\cref{lem:front12:reference}(i) with
$j=1$). By \cref{lem:tr:polarization}(a) with $\varepsilon=(f_1+h_1)r$,
\[
 \Pi^{(1)}_{E_0}+\Pi^{(1)}_{E_1}+(f_1+h_1)r\bigl(\Pi^{(1)}_{G_{00}}+\Pi^{(1)}_{G_{11}}\bigr)
 \succeq_{\vartheta_0}\Pi^{(1)}_{F_1}+\Pi^{(1)}_{F_2} .
\]
A diagonal copy preserves this order of cut matrices (\cref{lem:five-type}(e));
multiplying by the record weight $v_{\rm D}r^2$ gives
$E_{\rm add}+G\succeq_{\vartheta_0}F_{\rm fold}$ (\cref{rem:polarization-uses}(1)).
The reference is positive at $\vartheta_0$ (\cref{lem:front12:reference}), and
$\rho_{\Blk}-\mu_{\Blk}$, which is $E_{\rm add}+G-F_{\rm fold}$ divided by
$\widetilde\mu_{\Blk}(1)$, is a real combination of the three positive records,
of variation at most their total mass.
\end{proof}

So the first stage costs only second order in $r$ (the removed fold copies have mass $2v_{\rm D}r^2$): the nonaccepting
one-layer fold copies are traded for matching records at the one
cut $\vartheta_0$ where positivity is needed, and the row errors
are paid by merges of degree three.
Since only actual components of the reference are subtracted, $\rho$ and
$\lambda$ stay nonnegative. The tensor step, $Y_\rho\ge Y_\mu$, is Step~1 of
\cref{lem:record:first-stage}.

\paragraph{The folded matching records.}
For $s\in\{0,1\}$ the folded matching record of type $E_s$ in
$E_{\rm fold}+E_{\rm add}$ has mass $(v_{\rm S}+v_{\rm D})r^2$: the one-layer law
$\Pi^{(1)}_{E_s}$, uniform on the halves, placed on the reserved diagonal. By
\cref{cert:record:middle-geometry}, at its first-coordinate cut
($\vartheta_1$ for $E_0$, $\vartheta_2$ for $E_1$) the label of object~$2$ lies
on both sides of the cut (it is \emph{shared}), the labels of objects $0$
and $1$ lie on one side each (they are \emph{private}), and the six queries
need exactly the edge $y_0\to y_1$. At one prime (the block statement then follows by tensoring), given the shared label $c$, the private
pair is uniform on the ordered distinct pairs of $V_c=\Fp\setminus\{c\}$.
Write
$V_c^\pm=\{a\in V_c:\chi(a-c)=\pm1\}$, each of size $(p-1)/2$, let
$\mathbf M_c$ be the $0/1$ matrix of the relation $a\to a'$ on $V_c$, and
$\mathbf M_c^{+-}$ its restriction to $V_c^+\times V_c^-$, extended by zero.

\begin{lemma}[One-stage replacement of the folded matching records]\label{lem:record:one-stage-folded-e}
Let $p\ge11$ with $p\equiv3\pmod4$, $k=(p-3)/4$,
$\kappa^{\rm fold}_p=(1+3r)r/(1-3q)=(\sqrt p+3)/(p-3)$ and
\[
 k_{\rm new}(p)=\frac{1+5q}{(1-3q)(1+q)}+\frac{qr}{(1-3q)(1-2q)},\qquad
 s_{\rm fold}(p)=\frac{(1+r)^2+4q}{(1-3q)^2}+\frac q{(1-2q)^2}.
\]
\begin{enumerate}[label=(\alph*)]
\item The four blocks of $\mathbf M_c$, from $V_c^+$ to $V_c^+$, from $V_c^+$
 to $V_c^-$, from $V_c^-$ to $V_c^+$ and from $V_c^-$ to $V_c^-$, have equal
 row and column degrees, namely $k$, $k+1$, $k$ and $k$.
\item The matrix $\mathcal K_+=\mathbf M_c/(2k)-\mathbf M_c^{+-}/(2k(k+1))$ is
 doubly stochastic, supported on the pairs $a\to a'$, with entries $1/(2k)$
 or $1/(2(k+1))$ there; from each $a$ it chooses each destination sector
 with probability $1/2$ and, within a sector, is uniform on the out-neighbours of $a$ in that sector.
\item Let $\mathcal K_0=(\Jall-I)/(p-2)$ on $V_c$, the conditional coupling of
 the record being replaced, and $\Delta=\mathcal K_+-\mathcal K_0$. On $L^2$ of
 the uniform law of $V_c$ both $\mathcal K_+$ and $\mathcal K_0$ fix constants
 and preserve mean-zero functions; there
 $\|\mathcal K_+\|\le\kappa^{\rm fold}_p$ and $\|\mathcal K_0\|=1/(p-2)$; the
 symmetric part of $\Delta$ has norm exactly $1/(p-2)$, attained at the
 sector contrast $\one_{V_c^+}-\one_{V_c^-}$; and
 $\|\Delta\|\le k_{\rm new}(p)\,r$.
\item Let $\Blk$ be a block of such primes $p\ge X\ge2^{24}$ with least prime
 $b$, put $r_X=X^{-1/2}$, $q_X=1/X$, $\alpha_X=(1+3r_X)/(1-3q_X)$ and
 $\beta_X=1/(1-2q_X)$, write $k_{\rm new}(X)$ and $s_{\rm fold}(X)$ for the displayed formulas evaluated at $q_X,r_X$, and suppose
 $\alpha_X^2r_X+\beta_X^2q_Xr_X<k_{\rm new}(X)$. For every shared label
 $c=(c_p)_{p\in\Blk}$ the difference
 $\Delta_c=\bigotimes_p\mathcal K_+^{(c_p)}-\bigotimes_p\mathcal K_0^{(c_p)}$
 satisfies $\|\Delta_c\|\le k_{\rm new}(X)\,b^{-1/2}$ and
 $\|\operatorname{Sym}\Delta_c\|\le s_{\rm fold}(X)/b$.
\item Replace each folded matching record of type $E_s$ by the measure of
 the same mass under which the shared label is uniform and, given it, the
 private pair is distributed by $\bigotimes_p\mathcal K_+^{(c_p)}$, placed on
 the reserved diagonal; let $E_{\rm new}$ be the sum of the two replacements.
 Then $E_{\rm new}$ is a nonnegative reserved-diagonal measure of the mass of
 $E_{\rm fold}+E_{\rm add}$; it accepts the six queries; it is invariant under
 the groups $\mathcal G_p$ of the maps $z\mapsto\alpha z+\beta$ ($\alpha$ a
 nonzero square) acting independently at the primes of $\Blk$, so it has
 uniform one-slot marginals and its queried pair is uniform on the
 strict-square pairs; and at its own first-coordinate cut each of its
 records has the same actual conditional marginals as the record it
 replaces, differs from it by at most $k_{\rm new}(X)\,b^{-1/2}$ in sharp
 operator norm, uniformly in the shared label, its symmetric part by at
 most $s_{\rm fold}(X)/b$, and by at most twice its mass in variation. The coupling is direct: no label is resampled when the two private labels coincide.
\end{enumerate}
\end{lemma}

\begin{proof}
(a) We use the Paley counts of \cref{lem:paley-triple}(a): two distinct
labels have exactly $k$ common out-neighbours and $k$ common in-neighbours,
and $\#\{c':a\to c'\to b'\}=k+1-\one[a\to b']$ for $a\ne b'$. Now, for
$a\in V_c^+$, i.e.\ $c\to a$, the out-neighbours of $a$ in $V_c^+$ are the
common out-neighbours of $c$ and $a$ ($k$), and those in $V_c^-$ are the
labels on paths $a\to a'\to c$ ($k+1$, as $a\not\to c$); for $a\in V_c^-$ they
are the common out-neighbours of $a$ and $c$ ($k$) and the labels on paths
$a\to a'\to c$ ($k$, as $a\to c$). The column degrees are computed in the
same way from common in-neighbours and paths $c\to a\to a'$.
(b) By (a), the row of $a$ puts mass $1/2$ uniformly on each destination
sector. Since row and column degrees agree in each block, the column of $a'$
receives from each of the two source sectors its degree in that block
divided by twice that degree, in all $1/2+1/2=1$.
(c) Both matrices are doubly stochastic for the uniform law, which gives
the first claim. On $V_c$, $\mathbf M_c=(\Jall-I+\mathbf C_c)/2$, where
$\mathbf C_c$ is the compression of the $p\times p$ matrix $(\chi(a'-a))$,
whose norm is $\sqrt p$ (\cref{lem:square-shift}(c)); as $\chi(-1)=-1$,
$\mathbf C_c$ is skew-symmetric. On mean-zero $f$, $\Jall f=0$, so
$\|\mathbf M_cf\|\le(1+\sqrt p)\|f\|/2$. Each block of $\mathbf M_c^{+-}$ has
row and column sums $k+1$, so $\|\mathbf M_c^{+-}\|\le k+1$ by the Schur test.
Hence $\|\mathcal K_+f\|\le[(1+\sqrt p)/(4k)+1/(2k)]\|f\|=\kappa^{\rm fold}_p\|f\|$.
On constants $\Delta$ vanishes; on mean-zero functions
$\mathcal K_0=-I/(p-2)$, and, since $2k=(p-3)/2$ and $2k(k+1)=(p-3)(p+1)/8$,
\[
 \Delta f=\frac{\mathbf C_cf}{p-3}-\frac f{(p-3)(p-2)}
  -\frac{8\,\mathbf M_c^{+-}f}{(p-3)(p+1)} .
\]
Let $\mathbf h_c=\one_{V_c^+}-\one_{V_c^-}$, the sector contrast, and let
$\mathcal Z_c$ be the space of functions with mean zero on each of the two
sectors, so that the mean-zero functions are the orthogonal sum of the
multiples of $\mathbf h_c$ and $\mathcal Z_c$. By (b) and its proof, for every
$a'$ the column of $a'$ receives mass $1/2$ from each source sector, and every
row puts mass $1/2$ on each target sector. Hence $\mathcal K_+$ and
$\mathcal K_+^*$ map every mean-zero function to a function with zero sum on
each sector, so they map $\mathcal Z_c$ into itself, and
$\mathcal K_+\mathbf h_c=\mathcal K_+^*\mathbf h_c=0$; both summands reduce
$\Delta$, and $\Delta\mathbf h_c=\Delta^*\mathbf h_c=\mathbf h_c/(p-2)$. On
$\mathcal Z_c$ the all-ones part of the block $\mathbf M_c^{+-}$ vanishes and
the identity has no entries in it, so there $\mathbf M_c^{+-}$ is half the
$(V_c^+,V_c^-)$ block of $\mathbf C_c$: its norm is at most $\sqrt p/2$, and,
being off-diagonal for the splitting of $\mathcal Z_c$ by the sectors, its
compression to $\mathcal Z_c$ has symmetric and skew parts of norm at most
$\sqrt p/4$ each. As the compression of $\mathbf C_c$ is skew, the symmetric
part of $\Delta$ on $\mathcal Z_c$ has norm at most
\[
 \frac1{(p-3)(p-2)}+\frac{2\sqrt p}{(p-3)(p+1)}\le\frac1{p-2};
\]
after clearing the positive denominators the inequality reads
$4q+2r(1-2q)/(1+q)\le1$, whose left side decreases in $p$ and is below
$0.82$ at $p=11$. So the symmetric part of $\Delta$ has norm $1/(p-2)$,
attained at $\mathbf h_c$. Finally, on $\mathcal Z_c$,
\[
 \|\Delta\|\le\frac{\sqrt p}{p-3}+\frac1{(p-3)(p-2)}+\frac{4\sqrt p}{(p-3)(p+1)}
 =\frac{\sqrt p(p+5)}{(p-3)(p+1)}+\frac1{(p-3)(p-2)}=k_{\rm new}(p)\,r,
\]
which is larger than $1/(p-2)$, the norm of $\Delta$ on the multiples of
$\mathbf h_c$.
(d) Both couplings fix the constants on both sides, so they are
sector-preserving in the sense of \cref{lem:sector}, and $\Delta_c$ and its
symmetric part map every sector $H_S$ (of exact support $S\subseteq\Blk$)
into itself; their norms are the largest over the sectors
(\cref{lem:sector}(a)). On $H_\emptyset$, $\Delta_c=0$. On $H_{\{p\}}$,
$\Delta_c$ is the factor $\Delta$ at $p$, so by (c) its norm is at most
$k_{\rm new}(p)r_p$ and its symmetric part has norm $1/(p-2)$. The function
$k_{\rm new}$ is nondecreasing in $q$ and $r$ (its first term is
$(1+4q/(1+q))/(1-3q)$), so $k_{\rm new}(p)\le k_{\rm new}(X)$, while
$r_p\le b^{-1/2}$ and $1/(p-2)\le1/(b-2)\le\beta_X/b$. On $H_S$ with
$|S|\ge2$ the operator is
$\bigotimes_{p\in S}\mathcal K_+^\circ-\bigotimes_{p\in S}\mathcal K_0^\circ$,
the centred factors having norms $\|\mathcal K_+^\circ\|\le\kappa^{\rm fold}_p
\le\alpha_Xb^{-1/2}$ and $\|\mathcal K_0^\circ\|=1/(p-2)\le\beta_X/b$, all
below one; so its norm is at most
$\alpha_X^2/b+\beta_X^2/b^2\le(\alpha_X^2r_X+\beta_X^2q_Xr_X)b^{-1/2}
<k_{\rm new}(X)b^{-1/2}$. This proves the first bound. For the symmetric
part, the decomposition of (c) gives on the mean-zero functions at one prime
$\|\operatorname{Sym}\mathcal K_+^\circ\|\le1/(4k)+1/(4k)=2q/(1-3q)$ and
$\|\operatorname{Skew}\mathcal K_+^\circ\|\le(\sqrt p+1)/(4k)=r(1+r)/(1-3q)$,
since $\mathbf C_c$ is skew and the symmetric and skew parts of the
off-diagonal block $\mathbf M_c^{+-}$ have norm at most $(k+1)/2$. On a
two-prime sector $\{p,p'\}$, as $\mathcal K_0^\circ$ is symmetric,
\[
 \operatorname{Sym}\Delta_c=\operatorname{Sym}\mathcal K_+^\circ\otimes
 \operatorname{Sym}\mathcal K_+^\circ+\operatorname{Skew}\mathcal K_+^\circ\otimes
 \operatorname{Skew}\mathcal K_+^\circ-\mathcal K_0^\circ\otimes\mathcal K_0^\circ ,
\]
of norm at most
$[4q_pq_{p'}+r_pr_{p'}(1+r_p)(1+r_{p'})]/((1-3q_p)(1-3q_{p'}))
+q_pq_{p'}/((1-2q_p)(1-2q_{p'}))\le s_{\rm fold}(X)/b$, using
$r_pr_{p'}\le1/b$ and $q_pq_{p'}\le q_X/b$. On a sector with $|S|\ge3$ the
whole operator has norm at most $\alpha_X^3b^{-3/2}+\beta_X^3b^{-3}<1/b$, as
$X\ge2^{24}$. Since $\beta_X\le s_{\rm fold}(X)$ and $1\le s_{\rm fold}(X)$,
every sector gives at most $s_{\rm fold}(X)/b$.
(e) Positivity, mass and the reserved-diagonal placement are clear. The
support lies on $a_0\to a_1$, strictly, at every prime, so the queried
edge is a strict-square edge and the three labels are distinct. The maps
of $\mathcal G_p$ preserve the sectors and the relation $\to$, so the law
is invariant; as $\mathcal G_p$ is transitive on labels and on the pairs
$a\to a'$, the one-slot marginals and the law of the queried pair are
uniform. At the cut the two half marginals are the laws of
$(c,a_0)$ and $(c,a_1)$; by (b) they are uniform on distinct pairs, as for
the record being replaced. The record difference is the coupling
difference $\Delta_c$ of (d) conditional on the shared label, which gives the
sharp bound and the bound of its symmetric part, and two probability laws
differ by at most $2$ in variation. No conditioning on a coincidence is
made.
\end{proof}

Only one strict edge has to be imposed, because the six queries need only
the edge between the two private labels of a folded matching record; a
single sector-balanced coupling of the private pair, which keeps both
actual half marginals at the record's own cut, therefore replaces each
record, with no additional conditioning event and no resampling at coincidences. The sector balance also
keeps the difference nearly skew: at a prime its symmetric part has norm exactly
$1/(p-2)$, carried by the sector contrast, and on a block it is of order
$1/b$, while the sharp coefficient $k_{\rm new}$ is about $1+7q$. The
symmetric part gives each folded tag the symmetric price
$(v_{\rm S}+v_{\rm D})s_{\rm fold}(X)r^4q_{\Blk}^{2a}$ on a block, which the
orbit coupon of the positive envelopes uses (\cref{lem:orbit-coupon}). The
coefficients
$k_{\rm new}(X)$ and $s_{\rm fold}(X)$ at the $33$ piece starts and the
two-support gate of (d) are recorded in \cref{cert:record:one-stage}
(\cref{app:certificates}); the gate holds at every $X\ge2^{24}$, since
$k_{\rm new}(X)>1$ while its left side is below $2^{-11}$. The replaced records
are positive records whose normalizations are admissible: they involve only
the field labels, and their sharp bounds hold for every value of the shared
label, which is (O2) of \cref{def:admissible}.

\emph{(CP2)(b).} On a middle block, $\delta_{E,\Blk}=\sum_{s\in\{0,1\}}
(\kappa'_s-\kappa_s)$, where $\kappa_s$ is the folded matching record of type
$E_s$ in $E_{\rm fold}+E_{\rm add}$ and $\kappa'_s$ its replacement in $E_{\rm new}$:
the two have the same mass $(v_{\rm S}+v_{\rm D})r^2$, and by
\cref{lem:record:one-stage-folded-e}(e) the difference has a sharp bound
$k_{\rm new}(X)b^{-1/2}$ relative to $\kappa_s$ at the cut $\vartheta_{1+s}$,
in the sense of \cref{def:price}(b), with the shared label as the label of
that definition, and variation at most twice the mass. The difference of type
$E_s$ carries the tag $s$ of (CP2)(b), whose cut $\vartheta_{1+s}$ is the same
in every middle block. This is (CP2)(b) with two tags, with the sharp and
variation prices of \cref{lem:record:middle-prices}; the tag totals are
equal.

\subsection{The law $\Nbr$ dominating $\lambda$ at $\iota$}\label{subsec:middle-gram}

Hypothesis~(CP3) requires one nonnegative law $\Nbr$ that dominates
$\lambda$ at the reserved cut in a $2\times2$ block. The next lemma
builds $\Nbr$ from $\mu$ and $\lambda$ record by record and proves
that the block is positive semidefinite.

\begin{lemma}[The $2\times2$ block dominating $\lambda$ at $\iota$]\label{lem:record:gram}
Let the weights of every middle piece satisfy the thirteen inequalities of
\cref{lem:front12:reference}(iii) and the protection \eqref{eq:protection}.
\begin{enumerate}[label=(\roman*)]
\item $\mu\preceq_\iota\Nbrz$, and $\Nbrz-\mu=E_{\rm add}+G+H$ is a
 nonnegative reserved-diagonal measure.
\item $\Nbr=\Nbrz+\delta_E=\lambda+F_{\rm fold}+H$ is a nonnegative measure,
 $\lambda_{\Blk}$ is symmetric at $\iota$, and blockwise
 $\begin{pmatrix}\Nbr_{\Blk}&\lambda_{\Blk}\\ \lambda_{\Blk}&\Nbr_{\Blk}\end{pmatrix}
 \succeq_\iota0$; in particular $-\Nbr\preceq_\iota\lambda\preceq_\iota\Nbr$.
\end{enumerate}
\end{lemma}

\begin{proof}
(i) Each of $E_{\rm add}$, $G$ and $H$ is a positive measure fixed
samplewise by $\iota$, so its cut matrix is a nonnegative diagonal matrix
(\cref{lem:commoncut}), and so is that of their sum $\Nbrz-\mu$.

(ii) $\lambda$ is nonnegative: $F_{\rm fold}$ and $E_{\rm fold}$ are actual
components of the unnormalized reference, and $E_{\rm new}$ and $G$ are
nonnegative. Hence so is $\Nbr=\lambda+F_{\rm fold}+H$. The reference is
invariant under the layer swap (\cref{lem:front12:reference}) and all records
are reserved diagonal, so the cut matrices of $\lambda_{\Blk}$ and
$\Nbr_{\Blk}$ at $\iota$ are symmetric. As in (i),
$\Nbr-\lambda=F_{\rm fold}+H\succeq_\iota0$. Next,
$\Nbr+\lambda=2\mu-F_{\rm fold}-2E_{\rm fold}+H+2E_{\rm new}+2G$, where
$E_{\rm new}$ and $G$ are nonnegative and reserved diagonal; it remains to show
$2\mu-F_{\rm fold}-2E_{\rm fold}+H\succeq_\iota0$, and we do so summand by summand
in $2\widetilde\mu_{\Blk}$ (all records are divided by the same normalizer
$\widetilde\mu_{\Blk}(1)$). Each fold component $\Pi_{F_s}$, $s\in\{1,2\}$,
enters $2\widetilde\mu_{\Blk}$ with
$2w_{\rm D}r\,\Pi_{F_s}+2v_{\rm D}r^2\operatorname{Diag}_\iota(\Pi^{(1)}_{F_s})$,
and $-F_{\rm fold}$ removes one $v_{\rm D}r^2\operatorname{Diag}_\iota(\Pi^{(1)}_{F_s})$;
what remains is positive at $\iota$ by \cref{lem:front12:reference}(iv). Each
matching component $\Pi_{E_s}$ enters with
$2w_{\rm S}r\,\Pi_{E_s}+2v_{\rm S}r^2\operatorname{Diag}_\iota(\Pi^{(1)}_{E_s})$;
$-2E_{\rm fold}$ removes its folded part, and $H$ adds
$2w_{\rm S}r\cdot3/(b-3)\cdot\operatorname{Diag}_\iota(\Pi^{(1)}_{E_s})$. The
component $\Pi_{E_s}$ is uniform on two disjoint injective triples, and its
cut operator at $\iota$ relative to its actual half law, the uniform law on
injective triples, which is the half law of $\Pi^{(1)}_{E_s}$, is the
disjoint-triple operator at each prime; on the block its spectrum lies in
$[-3/(b-3),1]$ (\cref{sec:frontier:empty-triple-negative}(c)). Hence
$\Pi_{E_s}+\frac3{b-3}\operatorname{Diag}_\iota(\Pi^{(1)}_{E_s})\succeq_\iota0$.
Every other component (the root, the aligned fold $\Pi_{F_0}$, the merges and
the constants) keeps both its two-layer part and its folded part, which
are positive together at $\iota$ by \cref{lem:front12:reference}(ii) and the
third line $v_{\mathfrak t}\ge k_{\mathfrak t}w_{\mathfrak t}$ of
\cref{lem:front12:reference}(iii). So $2\mu-F_{\rm fold}-2E_{\rm fold}+H\succeq_\iota0$,
and $\Nbr_{\Blk}\pm\lambda_{\Blk}\succeq_\iota0$, which is the block statement
for the symmetric matrices: conjugation by the unitary
$2^{-1/2}\bigl(\begin{smallmatrix}1&1\\1&-1\end{smallmatrix}\bigr)$ turns the
block matrix into the diagonal matrix with blocks $\Nbr_{\Blk}+\lambda_{\Blk}$ and
$\Nbr_{\Blk}-\lambda_{\Blk}$ (\cref{def:cp3}). Outside the middle,
$\Nbr_{\Blk}=\lambda_{\Blk}=\mu_{\Blk}\succeq_\iota0$.
\end{proof}

With the price $u_{\Blk}$ of $\Nbr_{0,\Blk}-\mu_{\Blk}$ given by
\cref{lem:record:middle-prices}, the lemma is hypothesis (CP3) of
\cref{thm:comparison} (\cref{def:cp3}). The
energy bound that the bracket buys (every structured energy of a
fixed selected law with $\Nbr$ outside is at most $B+o(1)$ times its price times
$Y$, the part without record differences costing at most $e^{U_0}$ and the
rest at most $R_{\rm err}$) is \cref{lem:bracket-energy}, which assumes only
(CP1), (CP2)(b), (CP3) and (CP8).

\begin{remark}[What the bracket buys]\label{rem:middle-bracket}
The bracket is a Gram bracket at $\iota$, not a domination at every cut: the
comparison law built later need not be reflection positive, and the comparison
is made at the reserved cut only. Two features of the comparison rest on
this. First, $\lambda_{\Blk}$ is symmetric at $\iota$, so in the expansion of
$Y_\nu-Y_\lambda$ every tensor term with an odd number of skew factors is
skew-adjoint and has zero real quadratic form on the identical reserved
half-functions; this is why the reserved cost is $B(e^S\cosh A-1)$ and not
$B(e^{S+A}-1)$, where $S$ and $A$ are the total symmetric and skew prices of \cref{lem:record:record-expansion}. Second, the conductors
of the index supports are held fixed through every Cauchy--Schwarz step:
for complementary fold maps $f_+,f_-$ and every slot-weight vector $w=(w_v)$,
$\frac12(\sum_vw_{f_+(v)}+\sum_vw_{f_-(v)})=\sum_vw_v$, so although a later
chamber fold can give four occurrences of a conductor on one energy and
none on the other, their geometric mean has the original two; only after the
complementary energies are combined are the conductor powers averaged. The
conductor moment of the matching term is therefore exactly $R_I^6$, also in
the mixed-cut remainder of \cref{lem:bracket-energy}
(\cref{lem:record:folded-energy}(e)).

The bracket carries one copy of the one-layer folded copies of the fold records, so that
$\Nbrz-\mu=E_{\rm add}+G+H$, whose coefficient of $r^2$, namely $2v_{\rm D}$, is half
that of the first-stage records $F_{\rm fold}+E_{\rm add}$
(\cref{lem:record:middle-prices}). Its part $E_{\rm add}+H$ is invariant under
the chamber reflections and the group $\mathcal G_6$ they generate, as $\mu$
is, and the positive envelopes use $\mu+E_{\rm add}+H$ as their invariant
background (\cref{lem:oct24:invariant-background}). The price is the
protection \eqref{eq:protection}, binding in every piece
(\cref{rem:middle-weights}), and the diagonal copy of the fold
records available to the joint Gram bound of the matching
envelopes has coefficient $v_{\rm D}$ per side
(\cref{prop:record:matching}).
\end{remark}

\subsection{Prices $L_0$ and $B$ for $\lambda$ and $\Nbr$}
\label{subsec:middle-frame-prices}

\begin{lemma}[Middle prices]\label{lem:record:middle-prices}
Let $\Blk$ be a block of a middle piece starting at $X$, with least prime
$b\ge X$ and $r=b^{-1/2}$. Evaluate every local coefficient below at $X$.
Then the prices of $\Blk$ are at most $q_{\Blk}^{2a}$ times
\begin{align*}
 &\Nbrz-\mu:\ 2v_{\rm D}r^2+\Bigl(2v_{\rm D}(f_1+h_1)+\frac{12w_{\rm S}}{1-3/X}\Bigr)r^3,\\
 &F_{\rm fold}+E_{\rm add}+G:\ 4v_{\rm D}r^2+2v_{\rm D}(f_1+h_1)r^3,\\
 &\delta_E\ \text{(sharp)}:\ 2(v_{\rm S}+v_{\rm D})k_{\rm new}(X)\,r^3,\qquad
 \delta_E\ \text{(variation)}:\ 4(v_{\rm S}+v_{\rm D})r^2,
\end{align*}
where $k_{\rm new}$ is the coefficient of \cref{lem:record:one-stage-folded-e};
each of the two tags carries half of the last two prices.
\end{lemma}

\begin{proof}
Dividing by the normalizing mass $\widetilde\mu_{\Blk}(1)\ge1$ only decreases
each mass, so it suffices to bound unnormalized masses. The masses of
$E_{\rm add}$, $G$ and $H$ (\cref{def:record:frame}(b)) give the first price
and those of $F_{\rm fold}$, $E_{\rm add}$ and $G$ the second, with
$3/(b-3)=3r^2/(1-3q)\le3r^2/(1-3/X)$. The two prices of $\delta_E$ are those
of (CP2)(b) above: two record differences, each of a record of mass
$(v_{\rm S}+v_{\rm D})r^2$, with sharp bound $k_{\rm new}(X)r$ and variation
at most twice that mass, and no coincidence term; each tag consists of one of
them. The row constants $f_1,h_1$ and the coefficient $k_{\rm new}$ are valid
at every block of the piece when evaluated at $X$
(\cref{lem:record:local-domain}(iii)).
\end{proof}

Summing over the blocks of each piece with the charges $b_d(X,Y)$ of
\cref{lem:record:block-charge} for a piece $[X,Y)$, writing $b_2,b_3$ for $b_2(X,Y),b_3(X,Y)$ below, and over the $33$ pieces, with the weights and
row constants of each piece,
\begin{gather*}
 U_0\le\sum_{\rm pieces}\Bigl[2v_{\rm D}b_2+\Bigl(2v_{\rm D}(f_1+h_1)
  +\frac{12w_{\rm S}}{1-3/X}\Bigr)b_3\Bigr],\qquad
 P\le\sum_{\rm pieces}\bigl[4v_{\rm D}b_2+2v_{\rm D}(f_1+h_1)b_3\bigr],\\
 D\le\sum_{\rm pieces}2(v_{\rm S}+v_{\rm D})k_{\rm new}(X)\,b_3,\qquad
 V\le\sum_{\rm pieces}4(v_{\rm S}+v_{\rm D})\,b_2 ,
\end{gather*}
each tag carrying half of $D$ and of $V$. In the certificate arithmetic all coefficients have nonnegative
lower endpoints in every piece, so upper charges times upper coefficients are
upper bounds. The certificate evaluates these sums from
\path{research/push_2026_10_25/global/refined_L570_source.json},
with exact finite block histograms and outward analytic charges, at the anchor exponent
$a_0=\PushAnchor$ at which the middle prices are tabulated
(\cref{cert:record:middle-price}): the totals there (primed, denoting evaluation at $a_0$) have the upper ends
$U_0'\le0.8469081935391430$, $P'\le1.6936504653996420$,
$D'\le4.741704122563\cdot10^{-5}$ and $V'\le3.387253370148178$, and each of
the two tags has sharp total at most $D'/2$ and variation total at most
$V'/2$; the four sums of each piece are listed in
\cref{tab:record:middle-charges}. At the baseline exponent $a$
every middle block has $\log q_{\Blk}\le L=570$
(\cref{lem:record:local-domain}(ii)), so all these prices are at most
$t=e^{2(a-a_0)L}$ times their values at $a_0$ (\cref{lem:record:repricing}),
and \cref{cor:first-stage-repricing} with the two-tag series of
\cref{cor:two-cut-series} at $n=4$ gives $U_0<0.846908245753$,
$B<2.332628665480$, $R_{\rm err}<2.042558854365\cdot10^{-4}$ and
$L_0<4.763323725665\cdot10^{-4}$ (\cref{def:closing-quantities}(2)); the repriced
tag totals, about $2.37\cdot10^{-5}$ (sharp) and $1.694$ (variation), satisfy
the hypotheses of that corollary, since the sharp total is below
the variation total and below $6$, and their product is below $36$. For
comparison, the one-sharp bounds
$De^{P+V}$ and $De^{U_0+V}$ of \cref{rem:cp:tags} would be about
$7.7\cdot10^{-3}$ and $3.3\cdot10^{-3}$: in the tagged series a term whose
differences all carry one tag pays no variation price, and a term with both
tags pays the variation of the less frequent tag only. The prices of the
comparison records of the middle range are those of
\cref{lem:middle-comparison-prices}, assembled from
\cref{prop:signed-bounds,lem:oct24:matching-star,lem:record:exact-marginal-e,lem:merge-repair}
with the same block charges.

\paragraph{Consequences and what to read next.}
This section depends on the following statements, grouped by topic.
\Cref{lem:record:class-three}, \cref{lem:record:local-domain} and
\cref{lem:record:block-charge} (pairs, domain and charges);
\cref{lem:layer12-geometry} (the half graphs and the symmetries); the caps of
the root: \cref{lem:layer12-caps,imp:local:collision-gap,cert:positive-root-caps}, with \cref{lem:spec:cutkernel}(f)
and \cref{lem:spec:tensorcomp}(b); \cref{def:positive-root} and
\cref{lem:positive-root-properties,lem:imp:coherent39-law}; \cref{lem:old-f-law} (the cap
$k_{\rm D}$); \cref{prop:middle-negative} in the form of
\cref{cor:influence-schur}, \cref{lem:row-comparison}, \cref{lem:four-count},
\cref{lem:sector} and \cref{lem:majorant} (the rows, the block steps and the
evaluation at the piece start); \cref{lem:five-type} with
\cref{lem:slot-array}, \cref{lem:marginals}, \cref{lem:commoncut} and
\cref{lem:record:fold-identification} (the completion);
\cref{lem:tr:polarization} (the first stage); \cref{lem:paley-triple} and
\cref{lem:square-shift} (the one-stage repair);
\cref{sec:frontier:empty-triple-negative} (the coefficient of $H$);
\cref{def:cp2,def:cp3,cor:two-cut-series} (the tags and the series);
\cref{lem:record:repricing} and \cref{cor:first-stage-repricing} (prices at
$a$). Certificates:
\cref{cert:record:middle-price,cert:layer12-grid,cert:l20-grid,cert:positive-root-caps,cert:record:one-stage,cert:record:old-f-cap,cert:record:frame-data,cert:record:middle-geometry}.

\paragraph{Hypotheses verified in this section (\cref{tab:hypotheses}).}
(CP1) for the middle atoms, that is, (R1)--(R3) with (R2) at every cut of the
family (\cref{lem:front12:reference}); (CP2)(a) at $\vartheta_0$
(\cref{lem:middle-polarization}) and (CP2)(b) with two tags
(\cref{lem:record:one-stage-folded-e}, \cref{cert:record:one-stage}); (CP3)
(\cref{lem:record:gram}, under the protection \eqref{eq:protection}); the
prices $P$, $V$, $D$, $U_0$ and the tag totals
(\cref{lem:record:middle-prices,lem:record:block-charge,cert:record:middle-price},
\cref{lem:record:repricing}); and the admissibility clause (CP8)(a) for the
reference, the first-stage records and the bracket records. These are the
rows ``(CP1) middle'', ``(CP2)(a)'', ``(CP2)(b)'', ``(CP3)'' and ``$P$,
$V$, $D$, $U_0$, tags'' of \cref{tab:hypotheses}.
